\documentclass[aspectratio=169,xcolor=dvipsnames]{beamer}
\usetheme{SimpleDarkBlue}
\usepackage{hyperref}
\usepackage{graphicx} % Allows including images
\usepackage{booktabs} % Allows the use of \toprule, \midrule and \bottomrule in tables
\usepackage[spanish]{babel}
\usepackage[utf8]{inputenc}
\usepackage{array}
\usepackage{tabularx}
\usepackage{caption}
\usepackage{adjustbox}
\usepackage{amssymb}
\usepackage{tcolorbox}
\usepackage{amsmath}

\graphicspath{{./graficas}} 

\title{Estadística aplicada}
\subtitle{Conceptos previos: conjuntos, funciones, cálculo y probabilidad}

\author{Práxedis Jiménez Ruvalcaba}

\institute
{

}
\date{\today}

\tcbuselibrary{theorems}

% Custom colored theorem box
\newtcbtheorem{mytheo}{Teorema}%
{colback=orange!5!white,colframe=orange!75!black,fonttitle=\bfseries}{th}

\begin{document}

\begin{frame}[c]
    \begin{beamercolorbox}[sep=8pt,center,rounded=true,shadow=true]{title}
        \usebeamerfont{title}\inserttitle\par
        \ifx\insertsubtitle\@empty
        \else
          \vskip0.25em
          {\usebeamerfont{subtitle}\insertsubtitle\par}
        \fi 
    \end{beamercolorbox}

    \vspace{1cm}
    \begin{center}
        \usebeamerfont{author}\insertauthor\par
        \vspace{0.5cm}
        \usebeamerfont{institute}\insertinstitute\par
        \usebeamerfont{date}\insertdate
    \end{center}
    
\end{frame}

\section{Modelado de Regresión}
    \begin{frame}[fragile]{Regresión lineal}
        Si se tiene un vector de datos observables \textit{x} y uno de variables respuestas \textit{y}, entonces
        \begin{verbatim}
            y ~ x
            y ~ 0 + x
        \end{verbatim}
        Definen un modelo lineal simple, el primero con interessción y el segundo sin. Si se busca 
        obtener la logarítmica, entonces
        \begin{verbatim}
            log(y) ~ x
        \end{verbatim}
    \end{frame}

    \begin{frame}[fragile]{Otras regresiones}
        Para regresiones multivariadas, el operador + añade variables predictoras a la construcción de un modelo
        \begin{verbatim}
            y ~ x_1 + x_2 + ... + x_n
        \end{verbatim}
        Para regresiones curvas, podemos definir un modelo polinomial con
        \begin{verbatim}
            y ~ poly(x, grado)
        \end{verbatim}
    \end{frame}
\end{document}